% GENERATED FILE. Edit canonical lessons, then run npm run generate:books.
% canonical-source-hash: fnv1a64:5c940c1890fc895e
% canonical-lessons: HI-C157-taza, HI-C157-amir, HI-C157-garib, HI-C157-javan, HI-C157-burha

\chapter{Fresh, Rich, Poor, Young, Old}
\label{ch:hi-fresh-rich-poor-young-old}
\hlchaptermodality{157}

\begin{chapteropening}
\textbf{By the end of this chapter:} \emph{I can say fresh, rich, poor, young and old.}

Complete the last lesson of chapter 157: I can say fresh, rich, poor, young and old.
\end{chapteropening}

\section[\texorpdfstring{\hi{ताज़ा} (tāzā)}{ताज़ा (tāzā)}]{\hi{ताज़ा} (tāzā) — fresh}

\label{lesson:HI-C157-taza}



\begin{quote}
Before the new one: say the Hindi for salty, then the Hindi for spicy.
\end{quote}

\subsection*{You'll want to know: \hi{ताज़ा}}
\textbf{\hi{ताज़ा}} — \emph{tāzā} — \textquotedblleft{}fresh\textquotedblright{}.

\begin{cousinweb}
\textbf{\hi{ताज़ा}} (\emph{tāzā}) is a loan, usually traced to Persian \emph{tāza}, \textquotedblleft{}fresh\textquotedblright{}.
\end{cousinweb}

\begin{grammarlens}[title={What you've built}]
One more word for this chapter.
\end{grammarlens}

\subsection*{Your turn}
\begin{itemize}
\raggedright
  \item \emph{Say it:} \emph{tāzā}
  \item \emph{Say it:} \emph{tāzā}, once more
  \item \emph{Say it:} say \emph{tāzā}
  \item \emph{Recall:} say the Hindi for salty, then the Hindi for spicy, then say \emph{tāzā} again
\end{itemize}

\begin{tcolorbox}[breakable,colback=teal!4,colframe=teal!35!black,title={Before you move on}]
What does \hi{ताज़ा} mean? (\textquotedblleft{}Fresh\textquotedblright{}.) Say it once more.
\end{tcolorbox}
\section[\texorpdfstring{\hi{अमीर} (amīr)}{अमीर (amīr)}]{\hi{अमीर} (amīr) — rich}

\label{lesson:HI-C157-amir}



\begin{quote}
Before the new one: say the Hindi for spicy, then the Hindi for fresh.
\end{quote}

\subsection*{You'll want to know: \hi{अमीर}}
\textbf{\hi{अमीर}} — \emph{amīr} — \textquotedblleft{}rich\textquotedblright{}.

\begin{cousinweb}
\textbf{\hi{अमीर}} (\emph{amīr}) is a loan from Arabic \emph{amīr}, \textquotedblleft{}commander, prince\textquotedblright{}, which reached Hindi through Persian. Its Arabic root is \emph{ʾ-m-r} (Lane). Hindi uses it for someone rich.
\end{cousinweb}

\begin{grammarlens}[title={What you've built}]
One more word for this chapter.
\end{grammarlens}

\subsection*{Your turn}
\begin{itemize}
\raggedright
  \item \emph{Say it:} \emph{amīr}
  \item \emph{Say it:} \emph{amīr}, once more
  \item \emph{Say it:} say \emph{amīr}
  \item \emph{Recall:} say the Hindi for spicy, then the Hindi for fresh, then say \emph{amīr} again
\end{itemize}

\begin{tcolorbox}[breakable,colback=teal!4,colframe=teal!35!black,title={Before you move on}]
What does \hi{अमीर} mean? (\textquotedblleft{}Rich\textquotedblright{}.) Say it once more.
\end{tcolorbox}
\section[\texorpdfstring{\hi{ग़रीब} (ġarīb)}{ग़रीब (ġarīb)}]{\hi{ग़रीब} (ġarīb) — poor}

\label{lesson:HI-C157-garib}



\begin{quote}
Before the new one: say the Hindi for fresh, then the Hindi for rich.
\end{quote}

\subsection*{You'll want to know: \hi{ग़रीब}}
\textbf{\hi{ग़रीब}} — \emph{ġarīb} — \textquotedblleft{}poor\textquotedblright{}.

\begin{cousinweb}
\textbf{\hi{ग़रीब}} (\emph{ġarīb}) is a loan from Arabic \emph{gharīb}, \textquotedblleft{}a stranger\textquotedblright{}, which reached Hindi through Persian. Its Arabic root is \emph{gh-r-b} (Lane). Hindi uses it for \textquotedblleft{}poor\textquotedblright{}.
\end{cousinweb}

\begin{grammarlens}[title={What you've built}]
One more word for this chapter.
\end{grammarlens}

\subsection*{Your turn}
\begin{itemize}
\raggedright
  \item \emph{Say it:} \emph{ġarīb}
  \item \emph{Say it:} \emph{ġarīb}, once more
  \item \emph{Say it:} say \emph{ġarīb}
  \item \emph{Recall:} say the Hindi for fresh, then the Hindi for rich, then say \emph{ġarīb} again
\end{itemize}

\begin{tcolorbox}[breakable,colback=teal!4,colframe=teal!35!black,title={Before you move on}]
What does \hi{ग़रीब} mean? (\textquotedblleft{}Poor\textquotedblright{}.) Say it once more.
\end{tcolorbox}
\section[\texorpdfstring{\hi{जवान} (javān)}{जवान (javān)}]{\hi{जवान} (javān) — young}

\label{lesson:HI-C157-javan}



\begin{quote}
Before the new one: say the Hindi for rich, then the Hindi for poor.

Describe with words that never bend: \textbf{\hi{आसान}}, \textbf{\hi{मुश्किल}}, \textbf{\hi{सुंदर}}, \textbf{\hi{बदसूरत}}.
\end{quote}

\subsection*{You'll want to know: \hi{जवान}}
\textbf{\hi{जवान}} — \emph{javān} — \textquotedblleft{}young\textquotedblright{}.

\begin{cousinweb}
\textbf{\hi{जवान}} (\emph{javān}) is usually taken as a loan from Persian \emph{javān}, \textquotedblleft{}young\textquotedblright{}. That may not be the whole story: CDIAL lists a Hindi \emph{juwān}, \textquotedblleft{}a youth\textquotedblright{}, under \emph{*yuvāna}, \textquotedblleft{}young\textquotedblright{}, a form it builds on the Sanskrit \emph{yuvan}, \textquotedblleft{}young\textquotedblright{} (CDIAL 10510). So Hindi may have this word by an inherited road as well as a borrowed one.
\end{cousinweb}

\begin{grammarlens}[title={What you've built}]
One more word for this chapter.
\end{grammarlens}

\subsection*{Your turn}
\begin{itemize}
\raggedright
  \item \emph{Say it:} \emph{javān}
  \item \emph{Say it:} \emph{javān}, once more
  \item \emph{Say it:} say \emph{javān}
  \item \emph{Recall:} say the Hindi for rich, then the Hindi for poor, then say \emph{javān} again
\end{itemize}

\begin{tcolorbox}[breakable,colback=teal!4,colframe=teal!35!black,title={Before you move on}]
What does \hi{जवान} mean? (\textquotedblleft{}Young\textquotedblright{}.) Say it once more.
\end{tcolorbox}
\section[\texorpdfstring{\hi{बूढ़ा} (būṛhā)}{बूढ़ा (būṛhā)}]{\hi{बूढ़ा} (būṛhā) — old (of people)}

\label{lesson:HI-C157-burha}



\begin{quote}
Before the new one: say the Hindi for poor, then the Hindi for young.
\end{quote}

\subsection*{You'll want to know: \hi{बूढ़ा}}
\textbf{\hi{बूढ़ा}} — \emph{būṛhā} — \textquotedblleft{}old (of people)\textquotedblright{}.

\begin{cousinweb}
\textbf{\hi{बूढ़ा}} (\emph{būṛhā}) is a tadbhava, a word inherited and reshaped by sound change. It comes from \emph{*buḍḍha}, \textquotedblleft{}old\textquotedblright{}, a form CDIAL reconstructs from its descendants (CDIAL 9271). CDIAL notes that the initial \emph{b-} rules out direct descent from Sanskrit \emph{vṛddha}, \textquotedblleft{}grown old\textquotedblright{}, though that word may have influenced it.
\end{cousinweb}

\begin{grammarlens}[title={What you've built}]
One more word for this chapter.
\end{grammarlens}

\subsection*{Your turn}
\begin{itemize}
\raggedright
  \item \emph{Say it:} \emph{būṛhā}
  \item \emph{Say it:} \emph{būṛhā}, once more
  \item \emph{Say it:} say \emph{būṛhā}
  \item \emph{Recall:} say the Hindi for poor, then the Hindi for young, then say \emph{būṛhā} again
\end{itemize}

\begin{tcolorbox}[breakable,colback=teal!4,colframe=teal!35!black,title={Before you move on}]
What does \hi{बूढ़ा} mean? (\textquotedblleft{}Old (of people)\textquotedblright{}.) Say it once more.
\end{tcolorbox}
